%%%%%%%%%%%%%%%%%%%%%%%%%%%%%%%%%%%%%%%%%
% CV (Academic) -- revised structure + cleaned content
% Uses resume.cls (same as your current Overleaf setup)
%%%%%%%%%%%%%%%%%%%%%%%%%%%%%%%%%%%%%%%%%

\documentclass{resume}

% -------------------- Packages --------------------
\usepackage[hidelinks]{hyperref}
\usepackage{url}
\usepackage{multicol}
\usepackage[left=0.5in,top=0.6in,right=0.5in,bottom=0.6in]{geometry}

% Optional (keeps URLs from overflowing)
\sloppy

% -------------------- Header --------------------
\name{NAKUL WEWHARE}

% NOTE: resume.cls sometimes renders \address on a single line.
% The most reliable way is to use multiple \address commands (one per line).
\address{%
\href{mailto:nakul@princeton.edu}{nakul@princeton.edu} \;\textbullet\;
+1 (609) 255-3916}
\address{%
Princeton University, Dept.\ of Ecology \& Evolutionary Biology (EEB)}
\address{%
\href{https://scholar.google.com/citations?user=T08iq_UAAAAJ&hl=en}{Google Scholar} \;\textbullet\;
\href{https://github.com/Nakul-wewhare}{GitHub}}

\begin{document}

% -------------------- Research focus --------------------
\begin{rSection}{Research Focus}
Bioacoustics and animal communication, with an emphasis on (i) vocal labeling and matching of others in animals, (ii) the role of vocal communication in social coordination and cooperation, and (iii) vocal sequences and syntactic structure.
\end{rSection}

% -------------------- Education --------------------
\begin{rSection}{Education}
{\bf Princeton University} \hfill {\em Sep 2025 -- Present}\\
Ph.D. Student, Department of Ecology \& Evolutionary Biology (EEB)\\
Advisor: Dr.\ Gerald G.\ Carter

{\bf Indian Institute of Science Education and Research (IISER) Pune} \hfill {\em Aug 2020 -- May 2025}\\
BS--MS Dual Degree (Biology major) \hfill \\
M.S.\ thesis: \textit{Social network dynamics of budgerigars and their influence on vocal learning}\\
Advisor: Dr.\ Anand Krishnan
\end{rSection}

% -------------------- Publications --------------------
\begin{rSection}{Publications}
\item {\bf Wewhare, N.} \& Krishnan, A.\ (2024).
{\bf Individual-specific associations between warble song notes and body movements in budgerigar courtship displays.}
\textit{Biology Open}, 13(10): bio060497. \href{https://doi.org/10.1242/bio.060497}{https://doi.org/10.1242/bio.060497}\\
{\small Related profile: {\bf Wewhare, N.}\ (2024). \textit{First person -- Nakul Wewhare}.
\textit{Biology Open}, 13(10): bio061701. \href{https://doi.org/10.1242/bio.061701}{https://doi.org/10.1242/bio.061701}.}

\item Madabhushi, A.\ J., {\bf Wewhare, N.}, Binwal, P., Agarwal, V., \& Krishnan, A.\ (2023).
{\bf Higher-order dialectic variation and syntactic convergence in the complex warble song of budgerigars.}
\textit{Journal of Experimental Biology}, 226(20): jeb245678. \href{https://doi.org/10.1242/jeb.245678}{https://doi.org/10.1242/jeb.245678}

\item {\bf Wewhare, N.} \& Pandey, U.\ (2021).
{\bf First record of scavenging by a Checkered keelback, \textit{Fowlea piscator} (Schneider 1799) (Natriciidae) from Western Maharashtra, India.}
\textit{Reptiles \& Amphibians}, 28(2), 348--349. \href{https://doi.org/10.17161/randa.v28i2.15638}{https://doi.org/10.17161/randa.v28i2.15638}
\end{rSection}

\begin{rSection}{Preprints}

\item {\bf Wewhare, N.}, Burkart, J.\ M., \& Wierucka, K.\ (2026).
{\bf Sequence-level vocal convergence in common marmosets.}
\textit{bioRxiv}. \href{https://doi.org/10.64898/2026.03.20.713272}{https://doi.org/10.64898/2026.03.20.713272}

\item Sarkar, R., Pal, T.\ S., Maji, S., Nandi, S., Lahiri, A., Basumatary, A.\ K., GP, A., {\bf Wewhare, N.}, Jahan, N., Jakhete, S., Bandyopadhyay, S., Adhikary, R., Paul, T., Majumdar, S., \& Bhadra, A.\ (2025).
{\bf Sight, smell and more: What cues do free-ranging dogs use for decision-making while scavenging?}
\textit{arXiv}, 2512.00058. \href{https://doi.org/10.48550/arXiv.2512.00058}{https://doi.org/10.48550/arXiv.2512.00058}

\item Dhati, A.\ M., {\bf Wewhare, N.}, Devasthale, N.\ S., Sundar, S., Parulekar, N., \& Krishnan, A.\ (2025).
{\bf Vocal sequences of diverse parrots suggest an important role for note repetition.}
\textit{bioRxiv}. \href{https://doi.org/10.1101/2025.10.24.684295}{https://doi.org/10.1101/2025.10.24.684295}
\end{rSection}

% -------------------- Teaching --------------------
\newpage
\begin{rSection}{Teaching Experience}
\item {\bf Teaching Assistant, \textit{EEB 338 / LAS 351: Tropical Biology}} \hfill Prof.\ Gerald Carter\\
Assisted in an intensive field course in Panama on tropical ecology, biodiversity, biotic interactions, and tropical landscape evolution, including support for field, laboratory, and data-analysis components.
\end{rSection}

% -------------------- Talks & Posters --------------------
\begin{rSection}{Conferences and Workshops}
\item {\bf 20th International Society for Behavioral Ecology (ISBE)} \\
Turin, Italy. Jul 20--24, 2026.\\
{\it Contributed talk:} {\bf Evidence that vampire bats use vocal matching to address specific group members.}\\
Recipient of a travel award (\texteuro 1600).

\item {\bf Bioacoustics \& AI 101: Making AI Accessible to Everybody} \\
Organized by the Behavioural Ecology group, University of Copenhagen. Sep 2024 (online).\\
{\it Contributed talk:} {\bf Quantifying syntactic convergence in common marmoset vocal sequences.}
\href{https://www.youtube.com/watch?v=bZplC_TDW1M}{(Recording)}

\item {\bf 5th Conference of the Indian Society of Evolutionary Biologists (ISEB)} \\
Organized by the Indian Society of Evolutionary Biologists at Lucknow University. Oct 2024.\\
{\it Poster:} {\bf Individual-specific associations between warble song notes and body movements in budgerigar courtship displays.}

\item {\bf Winter School on Evolutionary Game Theory} \\
Organized by the Department of Mathematics, Shiv Nadar University, Delhi. Dec 2023.

\item {\bf 4th Conference of the Indian Society of Evolutionary Biologists (ISEB)} \\
Organized by the Indian Society of Evolutionary Biologists at Ahmedabad University. Dec 2022.
\end{rSection}

% -------------------- Research experience --------------------
\begin{rSection}{Research Experience}
\begin{rSubsection}{Master's thesis student}{\href{https://sites.google.com/view/eceb-lab}{ECEB Lab, EOBU, JNCSAR}}{Supervisor: Dr.\ Anand Krishnan}{May 2024 -- May 2025}
\item Designed and ran fission--fusion experiments in lab budgerigar colonies to test how social-network position predicts longitudinal change in vocal repertoires.
\item Built interaction- and association-based social networks from (i) manual behavioural observations and (ii) automated video capture using a multi-camera Raspberry Pi setup.
\item Linked network-derived predictors to multi-level acoustic features of warble song (note-, motif-, and sequence-level structure) to quantify social effects on vocal learning.
\end{rSubsection}

\begin{rSubsection}{Undergraduate researcher}{\href{https://www.dpz.eu/en/contact/staff/profile/person/kaja-wierucka.html}{Behavioral Ecology \& Sociobiology Unit, German Primate Center}}{Supervisor: Dr.\ Kaja Wierucka}{Dec 2023 -- Apr 2024}
\item Developed a sequence-comparison approach for marmoset vocal sequences inspired by bioinformatics-style alignment to quantify syntactic convergence.
\item Quantified vocal accommodation at multiple scales: sequence-level syntax and note-level spectral structure (MFCC-based features).
\item Found sequence-level convergence after pair formation in calls directed toward strangers, but not partners; individual-call acoustics did not change in either audience context.
\end{rSubsection}

\newpage
\begin{rSubsection}{Semester project student}{\href{https://raghavrajan.wixsite.com/rajanlab}{Rajan Lab, IISER Pune}}{PI: Dr.\ Raghav Rajan}{Aug 2023 -- Dec 2023}
\item Built an automated microcontroller-based training setup for zebra finches to sing in a delay-instructed paradigm, separating preparation from production.
\item Assisted with electrophysiological recordings across song-related nuclei and downstream analysis of preparatory activity.
\end{rSubsection}

\begin{rSubsection}{MITACS Globalink summer intern}{\href{https://sites.psych.ualberta.ca/animal-cognition-ualberta/}{Animal Cognition Research Group, University of Alberta}}{PI: Dr.\ Lauren Guillette}{May 2023 -- Aug 2023}
\item Conducted behavioural experiments and analysis to quantify zebra finch activity and personality-linked information use during nest building.
\item Built a DeepLabCut-based video scoring pipeline; trained a general zebra finch model to support future experiments.
\item Implemented automated tracking workflows using AnimalTA for scalable behavioural quantification.
\end{rSubsection}

\begin{rSubsection}{Winter project student}{\href{https://knpsi.in/about.php}{Indore Zoo and ECEB Lab, IISER Bhopal}}{PI: Dr.\ Anand Krishnan}{Dec 2022 -- May 2023}
\item Recorded and analysed vocal sequences of parrots (\textit{Psittacula, Agapornis, Melopsittacus, Trichoglossus}) at the Indore Zoo.
\item Conducted cross-species comparative analyses of syntactic structure in complex vocal sequences.
\item Developed a semi-Markov model to describe sequential vocal structure across genera.
\end{rSubsection}

\begin{rSubsection}{Semester project student}{\href{https://sites.google.com/view/eceb-lab}{ECEB Lab, IISER Bhopal}}{PI: Dr.\ Anand Krishnan}{Aug 2022 -- Aug 2023}
\item Quantified relationships between budgerigar courtship displays and warble-song syntax, testing generality and individual variability.
\item Modelled courtship behaviour and vocal sequences using HMMs and deep sequence models (LSTMs; seq2seq).
\end{rSubsection}

\begin{rSubsection}{Summer intern}{\href{https://sites.google.com/view/eceb-lab}{Ecological, Community and Evolutionary Bioacoustics (ECEB Lab), IISER Bhopal}}{PI: Dr.\ Anand Krishnan}{Jun 2022 -- Aug 2022}
\item Analysed warble-song syntax to identify higher-order dialect differences across colonies.
\item Tested convergence in syntactic structure following colony fusion using longitudinal recordings.
\end{rSubsection}

\begin{rSubsection}{Undergraduate researcher}{\href{https://sites.google.com/view/doglabiiserkolkata/home}{Dog Lab, IISER Kolkata}}{PI: Dr.\ Anindita Bhadra}{Dec 2021 -- May 2022}
\item Recorded and analysed free-ranging dog acoustic interactions in the field (continuous sampling over multiple contexts).
\item Analysed vocalisations for context-dependent functional referencing; designed and ran playback experiments to quantify receiver discrimination.
\end{rSubsection}

\begin{rSubsection}{Summer intern}{\href{https://sites.google.com/view/doglabiiserkolkata/home}{Dog Lab, IISER Kolkata}}{PI: Dr.\ Anindita Bhadra}{Aug 2021 -- Dec 2021}
\item Designed and ran field experiments testing olfactory vs.\ visual cues during food scavenging and how human interactions shape decision-making.
\end{rSubsection}
\end{rSection}

% -------------------- Awards & fellowships --------------------
\newpage
\begin{rSection}{Awards \& Fellowships}
\item {\bf Kishore Vaigyanik Protsahan Yojana (KVPY) SX Fellowship} \hfill {\em 2020 -- 2025}\\
Government of India fellowship supporting research-oriented students during pre-PhD training.

\item {\bf MITACS Globalink Research Internship Fellowship} \hfill {\em Summer 2023}\\
Competitive 12-week research internship at a Canadian university under faculty supervision.
\par
\end{rSection}

% -------------------- Academic service --------------------
\begin{rSection}{Academic Service}
\item Co-reviewer (with Dr.\ Anand Krishnan), \textit{Proceedings of the Royal Society B} (2024).
\end{rSection}

% -------------------- Selected coursework (compressed) --------------------
\begin{rSection}{Selected Coursework}
{\bf Quantitative:} Bayesian inference; causal inference; network algorithms; parameter estimation and inverse theory; numerical computation; nonlinear dynamics; bioinformatics.\\
{\bf Biology:} Advanced evolution; Ecology I and II; animal behaviour; comparative physiology; developmental biology; Neuroscience I and II.\\
{\bf Mathematics:} Linear algebra; multivariate calculus; probability; partial differential equations; ordinary differential equations; graph theory.
\par
\end{rSection}

% -------------------- Technical skills (tight + defensible) --------------------
\begin{rSection}{Technical Skills}
\item {\bf Programming:} Python, R, MATLAB
\item {\bf Bioacoustics:} Raven Pro; warbleR; MFCC, spectral feature extraction; sequence/syntax analysis, DTW, UMAP
\item {\bf Quantitative:} GLMMs; Bayesian modelling; causal inference; network analysis; agent-based simulations
\item {\bf ML (applied):} HMMs; LSTM/seq2seq models; variational autoencoders (VAEs); image-based segmentation methods
\item {\bf Hardware/field systems:} Raspberry Pi multi-camera setups; Arduino/microcontrollers
\end{rSection}

% -------------------- Languages (kept) --------------------
\begin{rSection}{Languages}
English (TOEFL iBT 109/120), Hindi, Marathi
\end{rSection}

% -------------------- Extracurriculars (optional; shortened) --------------------
\begin{rSection}{Outreach \& Extracurricular }
\item Volunteer, Disha (teaching underprivileged children): mentored 5 students (ages 8--15) in math, science, and English (two years).
\item Photography Club Coordinator (2021--2022); Photographer for institute events (2022--2023).
\end{rSection}

\end{document}
